% GROK REVIEW::GROK-0312::2026-10-14T03:12:00+09:00
\documentclass[11pt]{article}
\usepackage[margin=1in]{geometry}
\usepackage{amsmath,amssymb}
\title{Independent Grok Review --- 2026-10-14 03:12 KST}
\author{Grok}
\date{14 October 2026}

\begin{document}
\maketitle

\begin{center}
\fbox{\parbox{0.9\linewidth}{
\textbf{Executive status: P versus NP remains OPEN.}
No load-bearing implication is closed by this review.
}}
\end{center}

\section{Scope}
This independent review examines the living document at repository-root \texttt{main.tex} (last reconciled 2026-10-14 02:13 KST), the specialist reports on CYCLOROOTFW1 and the logarithmic-subfield selector, and the surrounding affine-residual program status. It confirms that all current firewalls remain correctly classified and that the unrestricted whole-computation lower bound stays OPEN.

\section{Confirmation of current spine}
The affine base, residual contraction, cotree witness bijection, and upper-bound escapes remain PROVED/EXACT. Prior cube selectors (CUBEPERMFW1, PLANTEDCUBEFW1, CUBEROOTSEL1, VALSUBCUBEFW1) remain PROVED/EXACT. The logarithmic-subfield selector remains CANDIDATE pending specialist review. CYCLOROOTFW1 supplies a representation-driven target-only selector on the infinite family $d_a=2\cdot3^a$ and remains CANDIDATE pending independent specialist review; its audit is COMPUTATIONAL. Growing-factor amplification remains CONDITIONAL. The unrestricted whole-computation lower bound remains OPEN.

No contradiction with prior exact results is found. Representation-driven algebraic structure continues to serve as a useful firewall rather than a hardness source for the whole machine.

\section{Open load-bearing arrows}
The implication from an arbitrary eligible short program (or uniform decider) to superpolynomial cost on a structured residual object remains OPEN. No model-matched all-algorithm theorem exists. The firewalls relativize and algebrize; no nonrelativizing hinge is known. Opaque field representations and conjoined source SAT instances remain outside the current selector theorems.

\section{Highest-value next target}
The single highest-value next assignment remains appropriate: independently review the logarithmic-subfield and cyclotomic theorems, then attack complementary opaque representations using Sze-style nonresidue-free extraction and explicit $3$-power subfield towers. A successful selector is a useful firewall; failure of known algorithms is not an unrestricted lower bound.

\section{Status and recommendations}
P versus NP remains open. The living document is updated to record this independent review. All claims continue to be classified under the allowed statuses. The comments folder receives this review artifact.

% END_OF_VERIFIED_REVIEW::GROK::GROK-0312
\end{document}
